\documentclass[11pt]{article}

\usepackage[T1]{fontenc}
\usepackage[margin=1in]{geometry}
\usepackage{graphicx}
\usepackage{booktabs}
\usepackage{hyperref}

\title{Scientific Discovery Agents for Chem/Bio Hackathon Workflows}
\author{Hackathon Agents Team}
\date{\today}

\begin{document}

\maketitle

\begin{abstract}
% TODO: Summarize the scientific question, agent workflow, main result, and evidence limits.
\end{abstract}

\section{Introduction}
% TODO: Motivate the discovery challenge and explain why an agentic workflow is useful.

\section{System Overview}
% TODO: Describe the planner, chemist, tool executor, critic, writer, and mechanism workflow.

\section{Methods}
% TODO: Document deterministic tools, model routing, run modes, and validation procedures.

\subsection{Agent Workflow}
% TODO: Describe bounded iteration, state sharing, and stop criteria.

\subsection{Scientific Tools}
% TODO: Describe RDKit descriptors, calculation tools, literature/RAG tools, DFT wrappers, and reporting.

\section{Experiments}
% TODO: Define benchmark tasks, inputs, run settings, and evaluation criteria.

\section{Results}
% TODO: Report validated outputs, generated artifacts, and quantitative summaries.

\section{Discussion}
% TODO: Interpret results proportionally and distinguish hypotheses from demonstrated findings.

\section{Limitations}
% TODO: List missing controls, model/tool failure modes, reproducibility gaps, and deployment constraints.

\section{Conclusion}
% TODO: State the narrow conclusions supported by the evidence and next work.

\bibliographystyle{plain}
\bibliography{references}

\end{document}
